\documentclass{article}
\usepackage{neurips_2025}

\usepackage[utf8]{inputenc}
\usepackage[T1]{fontenc}
\usepackage{hyperref}
\usepackage{url}
\usepackage{booktabs}
\usepackage{amsfonts}
\usepackage{nicefrac}
\usepackage{microtype}
\usepackage{xcolor}
\usepackage{amsmath}
\usepackage{graphicx}

\title{MedVQA-Router: Bridging Binary Classification and Open-Ended Captioning for Medical Visual Question Answering}

\author{%
  % TODO (Member 1): fill in your name / email
  Member 1 Name \\
  Department of Computer Science \\
  University of Houston \\
  \texttt{member1@example.com} \\
  \And
  Phuong Ho \\
  Department of Computer Science \\
  University of Houston \\
  \texttt{dho8@cougarnet.uh.edu} \\
  \And
  % TODO (Member 3): fill in your name / email
  Member 3 Name \\
  Department of Computer Science \\
  University of Houston \\
  \texttt{member3@example.com} \\
  \And
  Sana Azhar \\
  Department of Computer Science \\
  University of Houston \\
  \texttt{sanaazhar59@gmail.com} \\
}

\begin{document}

\maketitle

% ============================================================
\begin{abstract}
Medical Visual Question Answering (Med-VQA) requires simultaneous understanding
of clinical images and natural-language questions. We address both the binary
(yes/no) and open-ended subsets of Med-VQA using the
\texttt{robailleo/medical-vision-llm-dataset}~\cite{robailleo2024}, which
contains 4,793 samples across X-ray, CT, MRI, and ultrasound modalities. Our
two-stage pipeline routes questions via a rule-based classifier, directing
binary questions to a multimodal fusion model (ResNet-18 + BioBERT) and
open-ended questions to GPT-5 Mini via the OpenAI API. For binary VQA we
compare text-only baselines (TF-IDF + LR/SVM), an image-only CNN (ResNet-18),
and fusion models. Our best binary classifier---ResNet-18 fused with fine-tuned
BioBERT---achieves \textbf{79.5\% accuracy} and \textbf{0.83 AUC-ROC},
outperforming single-modality baselines by over 7 percentage points.
% TODO (Member 4): update abstract once captioning results are finalised
\end{abstract}


% ============================================================
\section{Introduction}

Medical Visual Question Answering (VQA) has attracted significant research
interest as a tool for automated clinical decision support. A Med-VQA system
must jointly reason over a clinical image and a natural-language question to
produce a correct answer---either binary (``Is there a fracture?'') or
open-ended (``What abnormality is visible?''), requiring different modelling
strategies.

Medical datasets are limited in size due to the high cost of expert annotation
and privacy regulations. Clinical images exhibit significant visual diversity
across modalities, while questions use a highly specialized vocabulary. Standard
vision-language models pretrained on natural images often fail to capture subtle
pathological markers critical for accurate diagnosis.

We propose a two-stage pipeline that prioritizes task specialization: a routing
classifier separates binary from open-ended questions, followed by a multimodal
fusion architecture combining fine-tuned ResNet-18 with domain-specific BioBERT
embeddings. This approach achieves 79.5\% accuracy and 0.83 AUC-ROC on binary
classification, significantly outperforming general-purpose models like CLIP.


% ============================================================
\section{Background}
\label{sec:background}

\textbf{Visual and Medical VQA.}
Since the VQA v2 dataset~\cite{goyal2017making}, VQA research has evolved from
monolithic architectures to complex attention-based mechanisms. In the medical
domain, datasets such as VQA-Med~\cite{lau2018dataset} and
SLAKE~\cite{liu2021slake} benchmark performance on radiology and pathology
images, where accuracy is safety-critical and linguistic bias is a known
challenge.

\textbf{Transfer Learning and Biomedical Language Models.}
Due to the scarcity of annotated medical images, transfer learning from
ImageNet-pretrained ResNet architectures~\cite{he2016resnet} is standard
practice. BERT~\cite{devlin2018bert} has spawned specialized variants such as
BioBERT~\cite{lee2020biobert}, pretrained on PubMed abstracts and PubMed
Central full texts, making it significantly better suited to clinical
terminology than general-domain models.

\textbf{Vision-Language Models.}
CLIP~\cite{radford2021clip} aligns images and text in a shared latent space
via dual encoders, while BLIP~\cite{li2022blip} enables flexible generation
through a multimodal encoder-decoder mixture. GPT-5 Mini extends this to
large-scale instruction-following captioning via the OpenAI API.


% ============================================================
\section{Data}
\label{sec:data}

% TODO (Member 1): expand this section with your full EDA findings

\subsection{Dataset}
We use the \texttt{robailleo/medical-vision-llm-dataset}~\cite{robailleo2024}
from HuggingFace, a combined medical vision-language dataset comprising
\textbf{4,793 samples} drawn from two sources: ROCO (3,000 samples) and
VQA-RAD (1,793 samples). The official split yields \textbf{3,834 training} and
\textbf{959 validation} samples.

\textbf{Modality distribution:} X-ray dominates with 2,325 samples (48.5\%),
followed by CT (1,351; 28.2\%), unknown modality (812; 16.9\%), MRI (231;
4.8\%), ultrasound (70; 1.5\%), and microscopy/endoscopy (4; $<$0.1\%).
\textbf{Body-part distribution:} Unknown (2,856), chest (718), abdomen (468),
head (408), extremities (125), spine (119), and pelvis (99).

\subsection{Preprocessing and Split}
Images are resized to $224\times224$ and normalized using ImageNet statistics.
Grayscale images are converted to three-channel RGB; questions are tokenized to
a maximum of 128 tokens. After routing, the binary subset contains \textbf{940
samples} (750 train / 190 val), nearly balanced at 49.5\% / 50.5\% yes/no. The
open-ended subset contains the remaining samples.
% TODO (Member 1): confirm binary/open-ended routing counts on full 4,793-sample dataset


% ============================================================
\section{Methods}
\label{sec:methods}

Our system has three components: (1) a routing classifier, (2) a binary VQA
pipeline, and (3) an open-ended captioning pipeline.

\begin{figure}[h]
  \centering
  \includegraphics[width=0.92\linewidth]{figures/pipeline.png}
  \caption{End-to-end pipeline: the routing classifier directs inputs to the
           binary classification or open-ended captioning branch.}
  \label{fig:pipeline}
\end{figure}

\subsection{Routing Classifier}
\label{sec:routing}
A rule-based router in \texttt{src/utils.py} inspects each answer: ``yes'' or
``no'' answers are directed to binary classification; all others to open-ended
captioning. The routing is deterministic and error-free by construction,
operating on ground-truth answer labels across all 4,793 samples.

\subsection{Text-Only Baselines}
Questions are vectorized with TF-IDF (10,000 terms, \texttt{ngram\_range=(1,3)},
\texttt{sublinear\_tf=True}). We train Logistic Regression (L2,
\texttt{max\_iter=1000}) and SVM (linear kernel) on these features. TF-IDF +
SVM achieves 68.95\% accuracy and 0.75 AUC-ROC, confirming meaningful but
insufficient signal from text alone.

\subsection{Image-Only Baseline}
ResNet-18~\cite{he2016resnet} pretrained on ImageNet is fine-tuned with
\texttt{layer4} unfrozen and a new binary classification head, reusing
low-level ImageNet features while adapting high-level representations to the
medical domain.

\subsection{Multimodal Fusion Models}
\label{sec:fusion}

Our fusion model concatenates image and text features before a learned head:
\begin{itemize}
  \item \textbf{Image encoder}: ResNet-18 (\texttt{layer4} unfrozen), 512-dim
        after global average pooling.
  \item \textbf{Text encoder}: \\Frozen CLIP (512-dim)~\cite{radford2021clip} or
        BioBERT \texttt{dmis-lab/biobert-base-cased-v1.2}~\cite{lee2020biobert}
        (768-dim, mean-pooled last hidden state).
  \item \textbf{Fusion head}: LayerNorm on each stream $\rightarrow$
        concatenation $\rightarrow$ Linear(256)--ReLU--Dropout(0.5)--Linear(1).
\end{itemize}

The fused representation is:
\begin{equation}
  z_{\text{fused}} = \bigl[\,\text{LN}(z_{\text{img}}) \oplus \text{LN}(z_{\text{txt}})\,\bigr]
\end{equation}
\hspace{4.65cm} where $z_{\text{img}} \in \mathbb{R}^{512}$, $z_{\text{txt}} \in \mathbb{R}^{768}$.

BioBERT is chosen over CLIP because it is pretrained on PubMed abstracts and
PubMed Central texts, making it better suited to clinical terminology. We
additionally fine-tune BioBERT layers 10--11 at $5\times10^{-6}$ to adapt text
representations while avoiding catastrophic forgetting.

Models are trained with \texttt{BCEWithLogitsLoss} and AdamW with separate LR
groups: $10^{-4}$ (fusion head), $10^{-5}$ (ResNet \texttt{layer4}), $5\times
10^{-6}$ (BioBERT layers 10--11). Batch size 16, weight decay $10^{-4}$, up to
30 epochs, early stopping (patience 7), \texttt{ReduceLROnPlateau} (factor 0.5,
patience 3).

\subsection{Open-Ended Captioning: GPT-5 Mini}
\label{sec:vlm}
% TODO (Member 3): write this entire subsection
For open-ended questions we query GPT-5 Mini (\texttt{gpt-5-mini-2025-08-07})
via the OpenAI API, sending each image as base64 alongside the question prompt.
% TODO (Member 3): describe prompt engineering, cost settings, and any
%   fine-tuning or few-shot strategy used


% ============================================================
\section{Experiments}
\label{sec:experiments}

\subsection{Evaluation Metrics}
Binary VQA: accuracy, macro F1, and AUC-ROC on the 959-sample validation split.
Open-ended captioning: BLEU-4, ROUGE-L, and
BERTScore~\cite{zhang2020bertscore}.

\subsection{Binary VQA Results}

\begin{table}[h]
  \caption{Binary VQA model comparison on the validation set.}
  \label{tab:binary-results}
  \centering
  \begin{tabular}{lccc}
    \toprule
    Model & Accuracy & F1 & AUC-ROC \\
    \midrule
    TF-IDF + Logistic Regression             & 0.6684 & 0.6441 & 0.7073 \\
    TF-IDF + SVM                             & 0.6895 & 0.6704 & 0.7517 \\
    ResNet-18 (layer4 unfrozen)              & 0.7263 & 0.7500 & 0.7740 \\
    Fusion: ResNet-18 + CLIP                 & 0.7579 & 0.7723 & 0.8194 \\
    Fusion: ResNet-18 + BioBERT              & 0.7789 & 0.7835 & 0.8204 \\
    Fusion: ResNet-18 + BioBERT (fine-tuned) & \textbf{0.7947} & \textbf{0.8000} & \textbf{0.8281} \\
    \bottomrule
  \end{tabular}
\end{table}

Results show a clear progression from text-only baselines to multimodal fusion.
Moving from ResNet-18 alone (72.63\%) to CLIP fusion (+3.16 pp) confirms that
visual features alone are insufficient without textual grounding. Replacing CLIP
with BioBERT adds a further 2.1 pp, reflecting its stronger clinical vocabulary
alignment. Fine-tuning the top two BioBERT layers provides the final boost to
79.47\% accuracy and macro-F1 of 0.80.

\subsection{Open-Ended Captioning Results}
% TODO (Member 3): fill in BLEU/ROUGE/BERTScore numbers from GPT-5 Mini

\begin{table}[h]
  \caption{Open-ended captioning results (GPT-5 Mini) on the validation set.}
  \label{tab:captioning-results}
  \centering
  \begin{tabular}{lccc}
    \toprule
    Model & BLEU-4 & ROUGE-L & BERTScore \\
    \midrule
    GPT-5 Mini (zero-shot)   & -- & -- & -- \\
    GPT-5 Mini (few-shot)    & -- & -- & -- \\
    \bottomrule
  \end{tabular}
\end{table}

\subsection{Qualitative Analysis}

\begin{figure}[h]
  \centering
  \includegraphics[width=0.82\linewidth]{figures/roc_curve.png}
  \caption{ROC curves for binary VQA models. The fine-tuned BioBERT fusion
           (AUC\,=\,0.8281) consistently outperforms the vision-only ResNet-18
           baseline (AUC\,=\,0.7740).}
  \label{fig:roc}
\end{figure}

\subsection{Compute Resources}
Binary fusion experiments used NVIDIA T4 and A100 GPUs via Google Colab Pro,
requiring $\approx$45 minutes to converge over 30 epochs (batch size 16,
max sequence length 128). Open-ended evaluation used the OpenAI API.


% ============================================================
\section{Conclusion}
\label{sec:conclusion}

We presented MedVQA-Router, a two-stage pipeline for Medical VQA that routes
questions to a binary classification branch or an open-ended captioning branch.
Combining ResNet-18 with fine-tuned BioBERT achieves 79.5\% accuracy and 0.83
AUC-ROC, outperforming all single-modality baselines by over 7 pp. The
consistent advantage of BioBERT over CLIP confirms the value of domain-adapted
language models for medical text understanding. Future work includes
cross-attention fusion for finer-grained image-text interaction, end-to-end
multi-task routing, and evaluation on larger benchmarks such as VQA-RAD and
SLAKE~\cite{liu2021slake}.


% ============================================================
\begin{ack}
This project was completed as part of the COSC 4368 Fundamentals of Artificial
Intelligence course at the University of Houston. We thank Professor Sen Lin
and the course teaching assistants for their guidance throughout the semester.
\end{ack}


% ============================================================
\bibliographystyle{plain}
\bibliography{references}


% ============================================================
\appendix

\end{document}
